\section{Finite spectral and projection facts}\label{sec:finite}
This section derives the local estimates used later; a regional gap is not
an additional hypothesis.

\begin{lemma}[Actual Fourier domain and periodic gap]\label{lem:fourier}
The dimension in \eqref{eq:domain} is $d_0=(D-1)(V-1)$. Every element of
$\UU_0$ has mean zero. On every componentwise mean-zero field,
\[
 \En(f)\geq\frac{16}{n^2}\norm f^2,
 \qquad\text{while on every field }\En(f)\leq4D\norm f^2.
\]
\end{lemma}
\begin{proof}
Use the counting-orthonormal characters
$V^{-1/2}\exp(2\pi i\ell\cdot x/n)$ with
$\ell\in\{0,\ldots,n-1\}^D$. The forward difference multiplier is
$\xi_i(\ell)=\exp(2\pi i\ell_i/n)-1$. Its energy multiplier is
$4\sum_i\sin^2(\pi\ell_i/n)$, which is at most $4D$.
At a nonzero frequency at least one term is at least
$4\sin^2(\pi/n)\geq16/n^2$, by concavity of sine on $[0,\pi/2]$;
the case $n=2$ gives equality and retains the actual forward labels.

The curl symbol is $v\mapsto\xi\wedge v$. When $\xi\ne0$, its kernel is
$\C\xi$, so the adjoint image has complex dimension $D-1$; at zero frequency
that image is zero. Conjugate frequencies pair to give twice their complex
dimension as a real dimension. At self-conjugate frequencies each $\xi_i$ is
real ($0$ or $-2$), and the real adjoint image has dimension $D-1$ directly.
Summing gives $(D-1)(V-1)$, without an omitted period-two mode. The zero
frequency is absent, proving the mean-zero assertion. The upper energy bound
also follows from the counting isometry of every shift.
\end{proof}

\begin{lemma}[Finite counting Neumann basis]\label{lem:neumann}
On an $\ell$-vertex path, $\ell\geq2$, the counting-orthonormal vectors are
\[
 \phi_0(t)=\ell^{-1/2},\qquad
 \phi_j(t)=\sqrt{2/\ell}\cos\frac{\pi j(t+1/2)}\ell
 \quad(1\leq j<\ell),
\]
with eigenvalues $4\sin^2(\pi j/(2\ell))$. A singleton has only its constant
mode. On any product rectangle with side lengths at most $L$, the counting
variance is at most $L^2/4$ times its internal edge energy, also for vector
fields; singleton factors contribute zero variance.
\end{lemma}
\begin{proof}
For $\theta=\pi j/\ell$, interior second differences of
$\cos(\theta(t+1/2))$ multiply it by $2(1-\cos\theta)$. The virtual endpoint
values $f(-1)=f(0)$ and $f(\ell)=f(\ell-1)$ give precisely endpoint degree
one. For an integer $1\leq r<2\ell$,
\[
 \sum_{t=0}^{\ell-1}\cos\frac{\pi r(t+1/2)}\ell
 =\frac{\sin(\pi r)}{2\sin(\pi r/(2\ell))}=0.
\]
The denominator is nonzero in this range. Product-to-sum, with the $r=0$
sum equal to $\ell$, proves the stated orthogonality and normalization.
The needed sums $j+j'$ are at most $2\ell-2$. There are $\ell$ orthonormal
vectors, so this is a complete basis.

Tensor products diagonalize the product-path energy. Since
$\sin t\geq2t/\pi$ on $[0,\pi/2]$, every nonconstant product mode has
eigenvalue at least $4/L^2$. Removing the constant mode gives the variance
bound by expansion in this finite basis. If the entire rectangle is a
singleton there is no positive eigenvalue to invoke and the variance is zero.
Sum the scalar inequalities for vector fields. No dimension factor is lost.
\end{proof}

For an integer $j\geq0$ partition each coordinate chart into the nonempty
intervals
\[
 \{\lfloor rn/2^j\rfloor,\ldots,\lfloor(r+1)n/2^j\rfloor-1\},
 \qquad0\leq r<2^j,
\]
and take their products. Call the partition $\mathcal P_j$ and its counting
averaging projection $\Pi_j$. Every coarse endpoint is an even-index fine
endpoint, so the partitions are nested after empty intervals are omitted.
There are at most $2^{Dj}$ cells and their sides have length at most
$L_j=\lceil n/2^j\rceil$. In particular $\Pi_0$ is global averaging.

\begin{lemma}[Actual spatial pieces]\label{lem:pieces}
Fix $k\geq0$. The operators
\[
 \Delta_j=\Pi_{j+1}-\Pi_j\quad(0\leq j<k),\qquad
 \Delta_k=I-\Pi_k
\]
are mutually orthogonal projections with sum $I-\Pi_0$. For every field $f$,
\begin{equation}\label{eq:piecegap}
 \norm{\Delta_jf}^2\leq\frac{L_j^2}{4}\En(f)
 \quad(0\leq j\leq k).
\end{equation}
For $j<k$, $\Delta_j$ is a subprojection of $\Pi_{j+1}$; the tail is a
subprojection of the singleton observation $I$.
\end{lemma}
\begin{proof}
Nested averaging projections satisfy $\Pi_i\Pi_j=\Pi_{\min(i,j)}$.
Their differences and the displayed tail are therefore orthogonal projections
and telescope. Apply Lemma~\ref{lem:neumann} on every parent cell of
$\mathcal P_j$ and sum. Its internal edges are a subset of the original
labelled periodic edges, counted once. Other edges add nonnegative energy.
This gives $\norm{(I-\Pi_j)f}^2\leq L_j^2\En(f)/4$.
Each $\Delta_j$ is a subprojection of $I-\Pi_j$, proving
\eqref{eq:piecegap}, including the single tail. The subprojection statements
follow from the same identities. Wrap edges are never inserted into a
Neumann cell as fictitious interior edges.
\end{proof}

\begin{lemma}[Finite kernel-to-all-field completion]\label{lem:completion}
Let $L$ be a positive semidefinite operator on a finite real Hilbert space,
let $R$ be an orthogonal projection, and suppose
$\ip u{Lu}\geq(c+\delta)\norm u^2$ on $\ker R$, with $c,\delta>0$.
If $\norm{L-cI}\leq M$, then
\begin{equation}\label{eq:completion}
 \ip z{Lz}+\left(M+\frac{2M^2}{\delta}\right)\norm{Rz}^2
 \geq c\norm z^2\quad\text{for all }z.
\end{equation}
For \eqref{eq:energy} with $0<B(x)\leq HI$, one may take $M=4DH+c$.
\end{lemma}
\begin{proof}
Write $z=u+v$, $u=(I-R)z$, $v=Rz$, and $F=L-cI$.
Then $\ip u{Fu}\geq\delta\norm u^2$, $\ip v{Fv}\geq-M\norm v^2$, and
\[
 2|\ip u{Fv}|\leq2M\norm u\norm v
 \leq\frac\delta2\norm u^2+\frac{2M^2}\delta\norm v^2.
\]
After adding the displayed penalty a nonnegative $\delta\norm u^2/2$ remains.
The final operator bound follows from Lemma~\ref{lem:fourier}'s upper energy
bound and $B(x)\leq HI$. Nothing requires $R$ to preserve a later restricted
domain. Constants must still be controlled by $R$ or a kernel margin; they
are not silently discarded by this lemma.
\end{proof}

\begin{lemma}[Original-coordinate singleton isolation]\label{lem:isolation}
Given $A$ distinct points in the original $D$-dimensional chart, there is a
partition into at most $1+2DA$ nonempty integer rectangles in which every
specified point is a singleton.
\end{lemma}
\begin{proof}
Start with the whole chart. In the rectangle containing the next point,
cut before and after its first coordinate and retain the two possible side
rectangles. Continue only within the middle slice, through the remaining
coordinates. At most $2D$ side rectangles and the desired singleton replace
one rectangle, a net increase at most $2D$. Omit empty pieces. Disjoint
coverage follows from the first coordinate at which a point differs from
the anchor. Earlier singleton anchors lie in other rectangles and remain
unchanged; an already-singleton point costs zero. Iteration proves the count.
This uses no global tensor-grid propagation of cuts or point separation.
\end{proof}
